% 数学建模竞赛论文LaTeX模板
% 适用于CUMCM（国赛）和MCM/ICM（美赛）

\documentclass[12pt,a4paper]{article}

% ============ 宏包导入 ============
\usepackage[UTF8]{ctex}                % 中文支持
\usepackage{amsmath,amssymb,amsfonts}  % 数学公式
\usepackage{graphicx}                  % 图片插入
\usepackage{float}                     % 浮动体控制
\usepackage{booktabs}                  % 三线表
\usepackage{geometry}                  % 页面设置
\usepackage{fancyhdr}                  % 页眉页脚
\usepackage{hyperref}                  % 超链接
\usepackage{listings}                  % 代码高亮
\usepackage{natbib}                    % 参考文献

% ============ 页面设置 ============
\geometry{top=2.5cm, bottom=2.5cm, left=2.5cm, right=2.5cm}

% ============ 页眉页脚 ============
\pagestyle{fancy}
\fancyhf{}
\fancyhead[C]{\small 数学建模竞赛论文}
\fancyfoot[C]{\thepage}

% ============ 代码高亮 ============
\lstset{
    basicstyle=\small\ttfamily,
    keywordstyle=\color{blue},
    commentstyle=\color{green!60!black},
    numbers=left,
    numberstyle=\tiny\color{gray},
    frame=single,
    breaklines=true,
}

\begin{document}

\title{\textbf{论文标题}}
\author{队员1 \quad 队员2 \quad 队员3 \\ \small 学校名称}
\date{\today}
\maketitle

% ============ 摘要 ============
\begin{abstract}
本文研究了\textbf{[问题名称]}问题。

\textbf{针对问题一}，建立了\textbf{[模型名]}，采用\textbf{[方法]}求解，得到\textbf{[结果]}。

\textbf{针对问题二}，构建了\textbf{[模型名]}，得到\textbf{[结果]}。

最后进行了灵敏度分析，结果表明\textbf{[评价]}。

\bigskip
\noindent\textbf{关键词：} 关键词1；关键词2；关键词3
\end{abstract}

\tableofcontents
\newpage

% ============ 一、问题重述 ============
\section{问题重述}
\subsection{问题背景}
[用自己的语言描述]

\subsection{问题要求}
\begin{enumerate}
    \item 问题一：[要求]
    \item 问题二：[要求]
\end{enumerate}

% ============ 二、问题分析 ============
\section{问题分析}
\subsection{问题一分析}
[分析问题本质，确定建模方向]

% ============ 三、模型假设 ============
\section{模型假设}
\begin{enumerate}
    \item 假设1：[内容]（合理性：[说明]）
    \item 假设2：[内容]
\end{enumerate}

% ============ 四、符号说明 ============
\section{符号说明}
\begin{table}[H]
\centering
\begin{tabular}{cll}
\toprule
\textbf{符号} & \textbf{含义} & \textbf{单位} \\
\midrule
$x$ & 变量 & 无 \\
$Z$ & 目标函数值 & 元 \\
\bottomrule
\end{tabular}
\end{table}

% ============ 五、模型建立与求解 ============
\section{模型建立与求解}
\subsection{问题一模型}
\subsubsection{模型建立}
\begin{equation}
\min Z = \sum_{i=1}^{n} c_i x_i
\end{equation}

\subsubsection{模型求解}
[求解步骤]

\subsubsection{结果分析}
\begin{figure}[H]
\centering
\includegraphics[width=0.8\textwidth]{figures/result.png}
\caption{结果可视化}
\end{figure}

% ============ 六、灵敏度分析 ============
\section{灵敏度分析}
[必须做！这是评奖的重要依据]

% ============ 七、模型评价 ============
\section{模型评价与推广}
\subsection{模型优点}
\subsection{模型缺点}
\subsection{改进方向}

% ============ 参考文献 ============
\begin{thebibliography}{99}
\bibitem{ref1} 作者. 标题[J]. 期刊, 年份.
\end{thebibliography}

% ============ 附录 ============
\appendix
\section{附录：代码}
\begin{lstlisting}[language=Python, caption={求解代码}]
import numpy as np
from scipy.optimize import linprog

c = [2, 3]
A = [[1, 1], [2, 1]]
b = [4, 5]
result = linprog(c, A_ub=A, b_ub=b)
print(f"最优解: {result.x}")
\end{lstlisting}

\end{document}
